\section{Background and Related Work}
\label{sec:background-related-work}

The literature relevant to this study spans cross-shard transaction protocols,
formal analysis of distributed and blockchain systems, implementation-model
validation, and reproducibility of verification artifacts. These areas contain
mature techniques used individually in our workflow, so we position the
contribution in their integration and controlled evaluation for one executable
cross-shard protocol.

\subsection{Sharded ledgers and cross-shard commit}
\label{subsec:related-cross-shard}

Sharded ledgers have long treated atomic cross-shard processing as a central
systems problem. Chainspace introduced S-BAC, a Byzantine distributed atomic
commit protocol for transactions whose objects span multiple shards, together
with correctness arguments and an executable prototype
evaluation~\cite{albassam2018chainspace}. OmniLedger introduced Atomix, a
client-driven atomic cross-shard transaction protocol that coordinates input
and output shards using proofs of acceptance or
rejection~\cite{kokoriskogias2018omniledger}. These systems establish that
atomic cross-shard commit mechanisms and protocol-level correctness arguments
predate the present work.

More recent systems focus on reducing conflicts, latency, or coordination
overhead. Prophet uses Byzantine-tolerant deterministic ordering to reduce
cross-shard conflicts and aborts and provides arguments for determinism and
serializability~\cite{hong2023prophet}. CSLAP builds leader accountability on
top of two-phase atomic commit and analyzes consistency, liveness, and
communication behavior~\cite{du2024cslap}. LightCross combines TEE-assisted
execution with a lightweight cross-shard commit protocol and presents
atomicity, serializability, and liveness arguments for its
design~\cite{qi2024lightcross}.

The cited cross-shard systems already provide correctness evidence through
protocol design, analytical or theorem-based arguments, and system-level
evaluation using metrics such as throughput and latency. Those metrics
characterize implementation behavior and should not be conflated with the
resource cost of formal verification. In contrast, RQ4 characterizes bounded
model-checking cost under a predefined experimental protocol. Mutation-based
property validation and implementation-model trace conformance are also not
reported as evaluation dimensions in the cited cross-shard works.

\subsection{Formal analysis of distributed and blockchain protocols}
\label{subsec:related-formal-analysis}

Formal modeling and model checking have been applied directly to distributed
transaction protocols. Schultz and Demirbas model MongoDB distributed
transactions in TLA+ and use TLC to analyze protocol properties. Their modular
treatment additionally generates tests from a formal storage model and
executes them against WiredTiger, providing a direct connection between the
model and a production implementation interface~\cite{schultz2025mongodb}.
This work is particularly close to our study because it combines cross-shard
distributed transactions, bounded formal exploration, and implementation-facing
validation.

Mutation-based assessment of formal specifications also predates this work.
Cordy et al. propose mutation model checking to assess and strengthen
specifications by checking whether deliberately mutated models are rejected by
the property set~\cite{cordy2023mutation}. Their work establishes
mutation-based specification assessment as a prior methodological approach.
Our RQ2 uses a smaller, protocol-targeted catalogue of scientific mutants and
interprets its mutation score only as sensitivity to those predefined defect
classes.

The use of multiple formal encodings and deliberate defects is likewise not
unique to this study. Konnov et al. translate the executable Ethereum 3SF
specification into TLA+, analyze it with Apalache, and use alternative SMT and
Alloy encodings for cross-validation~\cite{konnov2025ethereum}. Their technical
report also evaluates time, memory, timeouts, and deliberately injected bugs.
Consequently, neither the combination of TLA+ and Alloy nor the combination of
multi-formalism verification, defect injection, and verification-cost
measurements should be treated as an isolated novelty claim in our manuscript.

\subsection{Specification validation and implementation-model conformance}
\label{subsec:related-conformance}

A substantial body of work addresses the gap between verified models and
executable implementations. Cirstea et al. reduce validation of
distributed-program traces against TLA+ specifications to constrained model
checking with TLC and provide an API for instrumenting Java
programs~\cite{cirstea2025traces}. Their framework permits partially observed
traces and makes explicit the tradeoff between instrumentation detail and the
search required to determine whether an observed execution is compatible with
the specification.

TraceLink further automates the relation between distributed implementations
and verified designs. It links instrumented executions to TLA+ models,
evaluates trace-validation strategies, reports previously undetected
implementation defects, and publishes an experimental artifact with partial
external reproduction~\cite{hackett2025tracelink}. Smart Casual Verification
of CCF applies a related approach to a production C++ codebase by integrating
TLA+ model checking, simulation, testing, and trace validation into continuous
integration~\cite{howard2025ccf}. These studies show implementation-model trace
validation across multiple distributed-systems settings.

Other approaches generate or guide implementation checks from formal models.
Choreographic PlusCal projects distributed protocol descriptions to TLA+ and
generates conformance monitors, including evaluation on two-phase
commit~\cite{foo2023choreographic}. Mocket uses a verified TLA+ state space to
generate systematic tests, drive Java implementations through selected
behaviors, and compare runtime states with formal
counterparts~\cite{wang2023mocket}. iMocket restricts this strategy to
state-space regions affected by incremental changes~\cite{gao2025imocket}.
Beyond TLA+, the Amazon S3 ShardStore study combines executable reference
models, property-based testing, failure injection, and stateless model
checking~\cite{bornholt2021s3}.

These results constrain RQ3. We do not present trace validation or
implementation-model checking as a new general method. RQ3 evaluates bounded
implementation-model trace conformance for a predefined cross-shard catalogue
containing valid and deliberately corrupted traces with expected diagnostic
agreement. Its role is one implementation-facing layer in the evidence chain
used for RQ1, RQ2, and RQ4.

\subsection{Positioning of this work}
\label{subsec:related-positioning}

Within the comparator set for which full text was verified in our structured
review, we did not identify a cross-shard study that reports all of the
following within one frozen experimental protocol: bounded property
verification using complementary TLA+ and Alloy models, targeted mutation-based
property validation, bounded implementation-model trace conformance over both
valid and deliberately corrupted traces, explicit verification-cost
characterization, and independent reproduction of the analytical artifact.
This statement is deliberately scoped to the reviewed comparator set and is
not a priority claim about the complete literature.

Accordingly, DTL-Lab is positioned as an integrated and reproducible evidence
workflow for an executable cross-shard protocol. It connects bounded formal
verification, mutation sensitivity, implementation-model conformance, and
verification cost while preserving each layer's limitations, without claiming
the individual techniques as novel.
